\BourbakiExerciseGroup

\begin{enumerate}
\def\labelenumi{\arabic{enumi})}
\item
设 \(X_{1}, X_{2}\) 是两个有向集，f 是定义在 \(X_{1} \times X_{2}\) 上的实值函数，使得：对每个 \(x_{1} \in X_{1}\) ，映射 \(x_{2} \to f(x_{1}, x_{2})\) 在 \(X_{2}\) 上递增；对每个 \(x_{2} \in X_{2}\) ，映射 \(x_{1} \rightarrow f(x_{1}, x_{2})\) 在 \(X_{1}\) 上递增。试证：f 在 \(\overline{R}\) 中关于 \(X_{1}\) 与 \(X_{2}\) 的截面滤子之积有极限，并且该极限是 f 的上确界。
\item
  \begin{enumerate}
  \def\labelenumii{\alph{enumii})}
  \tightlist
  \item
设 \(f\) 是定义在集合 \(\mathbf{X}\) 上的实值函数，\(\mathbf{A}\) 是 \(\mathbf{X}\) 的一个非空子集，\(\varphi\) 是定义在 \(f(\mathbf{A})\) 上的递增实值函数。若 \(\varphi\) 在点 \(a = \sup_{x \in \mathbf{A}} f(x)\) 处连续，则 \(\sup_{x \in \mathbf{A}} \varphi(f(x)) = \varphi(\sup_{x \in \mathbf{A}} f(x))\) 。
  \end{enumerate}
\end{enumerate}

\begin{enumerate}
\def\labelenumi{\alph{enumi})}
\setcounter{enumi}{1}
\tightlist
\item
设 f 是定义在由滤子 F 滤子化的集合 X 上的实值函数，\(\varphi\) 是定义在 f 关于 F 的聚点所成之集的一个开邻域 V 上的实值函数。试证：若 \(\varphi\) 在 V 上递增且连续，则
\end{enumerate}

\[
\lim \sup _ {\mathfrak {F}} (\varphi \circ f) = \varphi (\lim \sup _ {\mathfrak {F}} f).
\]

\begin{enumerate}
\def\labelenumi{\arabic{enumi})}
\setcounter{enumi}{2}
\item
设 \(f\) 是定义在无穷集 \(\mathbf{X}\) 上的实值函数，\(\mathfrak{G}\) 是 \(\mathbf{X}\) 的有限子集的余集所成的滤子。试证：\(\limsup_{\mathfrak{G}} f\) 是所有这样的实数 \(x\) （集合 \(f([x, +\infty])\) 为无穷集）所成之集的上确界。
\item
设 \(f, g\) 是定义在集合 \(X\) 上的两个实值函数，该集合由滤子 \(\mathfrak{F}\) 滤子化。试证
\end{enumerate}

\[
\lim \sup _ {\mathfrak {F}} (\sup (f, g)) = \sup (\lim \sup _ {\mathfrak {F}} f, \lim \sup _ {\mathfrak {F}} g).
\]

试给出一个由滤子 F 滤子化的集合 X，以及定义在 X 上的实值函数的一个无穷族 \((f_{i})\) ，使得

\[
\sup _ {t} (\lim \sup _ {\mathfrak {F}} f _ {t}) <   \lim \sup _ {\mathfrak {F}} (\sup _ {t} f _ {t}).
\]

\begin{enumerate}
\def\labelenumi{\arabic{enumi})}
\setcounter{enumi}{4}
\tightlist
\item
设 \(f, g\) 是定义在滤子化集合上的两个实值函数。试证：若 \(\lim g\) 存在且 \(\geqslant 0\) ，则
\end{enumerate}

\[
\lim \sup f g = (\lim \sup f) (\lim g)
\]

只要两边都有定义。

\begin{enumerate}
\def\labelenumi{\arabic{enumi})}
\setcounter{enumi}{5}
\item
设 \(X_{1}\) （相应地 \(X_{2}\) ）是由滤子 \(F_{1}\) （相应地 \(F_{2}\) ）滤子化的集合，f 是定义在 \(X_{1} \times X_{2}\) 上的实值函数。试用例子证明：三个数 \(\limsup_{\mathfrak{F}_{i} \times \mathfrak{F}_{2}} f(x_{1}, x_{2})\) ， \(\limsup_{\mathfrak{F}_{i}} (\limsup_{\mathfrak{F}_{2}} f(x_{1}, x_{2}))\) ， \(\limsup_{\mathfrak{F}_{i}} (\limsup_{\mathfrak{F}_{i}} f(x_{1}, x_{2}))\) 一般是不同的。
\item
设 \(f\) 是定义在闭子集 \(A\) （ \(\overline{\mathbf{R}}\) 的）上的实值函数。用 \(\limsup_{x \to a, x \geqslant a} f(x)\) {[}相应地 \(\limsup_{x \to a, x \leqslant a} f(x)\) {]}表示 \(f(x)\) 当 \(x\) 趋于点 \(a \in A\) 且保持在 \(A \cap [a, +\infty]\) 中时的上极限。
\end{enumerate}


（相应地 A ∩ {[}---∞, a{]}）。试证：点 \(a \in A\) （满足 \(\limsup_{x \to a, x \geqslant a} f(x) \neq \limsup_{x \to a, x \leqslant a} f(x)\) ）所成之集是可数的。{[}证明：对每对有理数 p, q （满足 p \textless{} q ），考虑集 \(G_{p,q}\) ，即点 \(a \in A\) 中使

\[
\lim _ {x \rightarrow a, x \geqslant a} \sup f (x) \leqslant p <   q \leqslant \lim _ {x \rightarrow a, x \leqslant a} \sup f (x)
\]

成立的那些点所成之集；它是可数的（利用 § 2 习题 1）{]}。

\begin{enumerate}
\def\labelenumi{\arabic{enumi})}
\setcounter{enumi}{7}
\item
试由习题 7 推出：实值函数 \(f\) ，它定义在闭子集 \(\mathbf{A}\) （ \(\overline{\mathbf{R}}\) 的）上，且在 \(\mathbf{A}\) 上单调，则除 \(\mathbf{A}\) 的一个可数子集的点外，它在 \(\mathbf{A}\) 上连续。
\item
设 X 是具有可数基的拓扑空间，f 是定义在 X 上的实值函数。试证：点 \(x \in X\) （使 \(\lim_{x \to a, x \neq a} f(x)\) 存在且异于 \(f(a)\) ）所成之集是可数的（与习题 7 相同的方法）。
\item
设 X 是具有可数基 \((\mathbf{U}_{n})\) 的拓扑空间，f 是定义在 X 上的实值函数；称 f 在点 \(a \in X\) 处达到严格相对极大值，如果存在 a 的一个邻域 V，使得 \(f(x) < f(a)\) 对 \(x \in V\) 中每个异于 x = a 的点成立。试证：X 中使 f 达到严格相对极大值的点所成的集 M 是可数的。（考察这样的 \(U_{n}\) 所成之集：f 在 \(U_{n}\) 的某点处达到等于其在 \(U_{n}\) 中上确界的严格相对极大值，并证明存在这个集到 M 上的映射）。
\item
设 \((u_n)\) 是实数 \(>0\) 的序列，满足 \(\lim_{n\to \infty}u_n = 0\) 。试证：存在无穷多个指标 \(n\) ，使得 \(u_{n}\geqslant u_{m}\) 对一切 \(m\geqslant n\) 成立。
\item
设 \((u_n)\) 是实数 \(>0\) 的序列，满足
\end{enumerate}

\[
\lim _ {n \to \infty} \inf u _ {n} = 0.
\]

试证：存在无穷多个指标 \(n\) ，使得 \(u_{n} \leqslant u_{m}\) 对一切 \(m \leqslant n\) 成立。

\begin{enumerate}
\def\labelenumi{\arabic{enumi})}
\setcounter{enumi}{12}
\item
设 \((u_{n})\) 是有限实数的序列，\((\varepsilon_{n})\) 是数 \(\geqslant 0\) 的序列，满足 \(\lim_{n\to \infty}\varepsilon_n = 0\) ，且 \(u_{n + 1}\geqslant u_n - \varepsilon_n\) 对一切整数 \(n\geqslant 0\) 成立。令 \(a = \liminf_{n\to \infty}u_n\) ， \(b = \limsup_{n\to \infty}u_n\) ；试证：序列 \((u_{n})\) 的聚点所成之集是区间 \([a,b]\) 。
\item
设 \((r_n)\) 是有限实数 \(>0\) 的递增序列，满足 \(\lim r_n = +\infty\) 。对每个有限实数 \(r > 0\) ，用 \(\mathbf{N}(r)\) 表示最大指标 \(n\) ，它满足 \(r_n \leqslant r\) 。试证
\end{enumerate}

\[
\limsup _ {r \to \infty} \frac {\mathrm{N} (r)}{r} = \limsup _ {n \to \infty} \frac {n}{r _ {n}}, \quad \liminf _ {r \to \infty} \frac {\mathrm{N} (r)}{r} = \liminf _ {n \to \infty} \frac {n}{r _ {n}}.
\]

¶ 15) 设 \((x_n)\) 是有限实数的序列，\((p_n)\) 是有限实数 \(\geqslant 0\) 的序列，满足 \(\lim_{n\to \infty}\left(\sum_{i=0}^{n}p_i\right) = +\infty\) 。令 \(y_n = \left(\sum_{i=0}^{n}p_ix_i\right) / \left(\sum_{i=0}^{n}p_i\right)\) ，这里 \(n\) 取使 \(\sum_{i=1}^{\infty}p_i \neq 0\) 的那些值。试证 \(\liminf_{n\to \infty}x_n \leqslant \liminf_{n\to \infty}y_n \leqslant \limsup_{n\to \infty}y_n \leqslant \limsup_{n\to \infty}x_n\) 。

设 H 是序列 \((x_{n})\) 在 \(\overline{R}\) 中的聚点所成之集的任一非空子集。试证：序列 \((p_{n})\) （有限实数 \(\geqslant0\) 的）可以这样确定，使得 \(\lim_{n\to\infty}\left(\sum_{i=0}^{n}p_{i}\right)=+\infty\) ，并且使相应的序列 \((y_{n})\) 的聚点所成之集包含 H。{[}化归到 H 可数的情形，然后用归纳法定义 \((p_{n})\) ，取其各项为 o 或 1{]}。

由此推出：序列 \((x_{n})\) 在 \(\overline{\mathbf{R}}\) 中收敛，当且仅当序列 \((y_{n})\) 在 \(\overline{\mathbf{R}}\) 中收敛——对每个序列 \((p_n)\) （其项为数 \(\geqslant 0\) 且满足 \(\lim_{n\to \infty}\left(\sum_{i = 0}^{n}p_i\right) = +\infty .\) ）皆然。

\begin{enumerate}
\def\labelenumi{\arabic{enumi})}
\setcounter{enumi}{15}
\tightlist
\item
设 \(x_0, y_0\) 是满足 \(0 < y_0 \leqslant x_0\) 的两个实数。对 \(n > 0\) ，用下列关系式递归地定义两个序列 \((x_n), (y_n)\)
\end{enumerate}

\[
x _ {n + 1} = \left(x _ {n} + y _ {n}\right) / 2, \quad y _ {n + 1} = \sqrt {x _ {n} y _ {n}}.
\]

试证：两个序列 \((x_{n})\) ， \((y_{n})\) 趋于同一极限 \(\alpha\) （即 \(x_{0}, y_{0}\) 的``算术-几何平均''）；此外，存在数 \(a > 0\) 与 \(\gamma\) ，使得 \(0 < \gamma < 1\) 且 \(x_{n} - y_{n} \leqslant a\gamma^{2n}\) 对每个 \(n\) 成立{[}注意 \(x_{n+1} - y_{n+1} = (x_{n} - y_{n})^{2}/4(x_{n+1} + y_{n+1})\) {]}。

\begin{enumerate}
\def\labelenumi{\arabic{enumi})}
\setcounter{enumi}{16}
\tightlist
\item
设 \(g\) 是 \(]0, i]\) 到 \([-1, 1]\) 中的一个映射，满足
\end{enumerate}

\[
\lim _ {x \rightarrow 0, x > 0} g (x) = 0.
\]

试证：存在一个连续递增映射 \(g_{2}\) 与一个连续递减映射 \(g_{1}\) ，它们把 {[}0, 1{]} 映入 {[}-1, 1{]} ，且满足 \(g_{1}(0) = g_{2}(0) = 0\) 。

并且 \(g_{1}(x) \leqslant g(x) \leqslant g_{2}(x)\) 对 \(0 < x \leqslant 1\) 成立。{[}对每个整数 \(n > 0\) ，考察下确界 \(x_{n}\) ——诸数 \(x\) （满足 \(g(x) \geqslant 1 / n\) ）的下确界{]}。

\begin{enumerate}
\def\labelenumi{\arabic{enumi})}
\setcounter{enumi}{17}
\tightlist
\item
把第 1 节到第 6 节的定义与结果推广到这样的函数：其值取在于拓扑 \(\mathfrak{G}_{0}(\mathbf{X})\) 下为紧的全序集 X 中（参看 § 2，习题 6）。
\end{enumerate}

